% 强大数律的证明路线（教学示意，无数据）。
% 上排：Kolmogorov最大不等式 -> 方差可和的级数收敛；与Kronecker引理合并 -> 方差可和强大数律。
% 下排（上中下依次为方差、最终一致、均值）：截断 Y_n = X_n I{|X_n|<=n} 产生三条推导：
%   Borel--Cantelli（最终 X_n=Y_n）、截断项方差可和（sum Var(Y_n)/n^2 <= 2E|X_1|）、
%   控制收敛（E Y_n -> E X_1）；三者与方差可和强大数律合成有限一阶矩强大数律。
\begin{tikzpicture}[line width=.6pt,font=\fontsize{8}{10}\selectfont,>=Stealth,
  tool/.style={draw=FigureBlue,fill=FigureBlueFill,rounded corners=1.8pt,align=center,inner sep=3pt,minimum width=27mm,minimum height=9mm,text=FigureInk},
  stage/.style={draw=FigureGreen,fill=FigureGreenFill,rounded corners=1.8pt,align=center,inner sep=3pt,minimum width=31mm,minimum height=9mm,text=FigureInk},
  goal/.style={draw=FigureYellow,fill=FigureYellowFill,rounded corners=1.8pt,align=center,inner sep=3pt,minimum width=27mm,minimum height=9mm,text=FigureInk},
  flow/.style={->,draw=FigureStroke,line width=.7pt},
  lab/.style={font=\fontsize{7.5}{9}\selectfont,text=FigureMuted}]
 \node[tool] (max) at (0,0) {Kolmogorov\\最大不等式};
 \node[tool] (ser) at (42mm,0) {方差可和的\\独立级数收敛};
 \node[tool] (kro) at (84mm,14mm) {Kronecker引理};
 \node[goal] (sv) at (84mm,0) {方差可和条件下的\\强大数律};
 \node[stage] (tr) at (0,-30mm) {截断\\$Y_n=X_n\mathbb{I}_{\{|X_n|\leq n\}}$};
 \node[stage] (bc) at (42mm,-30mm) {Borel--Cantelli：\\最终$X_n=Y_n$};
 \node[stage] (va) at (42mm,-17mm) {$\sum_n\operatorname{Var}(Y_n)/n^2$\\$\leq2\,\mathbb{E}|X_1|$};
 \node[stage] (dc) at (42mm,-43mm) {控制收敛：\\$\mathbb{E}Y_n\to\mathbb{E}X_1$};
 \node[goal] (sl) at (122mm,-30mm) {有限一阶矩条件下的\\强大数律};
 \draw[flow] (max) -- (ser);
 \draw[flow] (ser) -- node[lab,above=1pt] {用于$W_n/n$} (sv);
 \draw[flow] (kro) -- (sv);
 \draw[flow] (tr) -- (bc);
 \draw[flow] (tr.east) -- (va.west);
 \draw[flow] (tr.east) -- (dc.west);
 \draw[flow] (va.east) -| node[lab,pos=.25,above=1pt] {对截断项应用} (sv.south);
 \draw[flow] (sv.east) -| (sl.north);
 \draw[flow] (bc.east) -- ([yshift=2mm]sl.west);
 \draw[flow] (dc.east) -- ([yshift=-2mm]sl.west);
\end{tikzpicture}
